\begin{thebibliography}{99}
\small
\bibitem{ref1} Hendrycks, D., Burns, C., Basart, S., et al. (2021). Measuring Massive Multitask Language Understanding (MMLU). ICLR 2021. arXiv:2009.03300. \url{https://arxiv.org/abs/2009.03300}
\bibitem{ref2} Liang, P., Bommasani, R., Lee, T., et al. (2023). Holistic Evaluation of Language Models (HELM). TMLR. arXiv:2211.09110. \url{https://arxiv.org/abs/2211.09110}
\bibitem{ref3} Liu, X., Yu, H., Zhang, H., et al. (2024). AgentBench: Evaluating LLMs as Agents. ICLR 2024. arXiv:2308.03688. \url{https://arxiv.org/abs/2308.03688}
\bibitem{ref4} Zhou, S., Xu, F. F., Zhu, H., et al. (2023). WebArena: A Realistic Web Environment for Building Autonomous Agents. arXiv:2307.13854. \url{https://arxiv.org/abs/2307.13854}
\bibitem{ref5} Atance, C. M., \& O'Neill, D. K. (2001). Episodic future thinking. Trends in Cognitive Sciences, 5(12), 533-539. \url{https://doi.org/10.1016/S1364-6613(00)01804-0}
\bibitem{ref6} Suddendorf, T., \& Corballis, M. C. (2007). The evolution of foresight: What is mental time travel, and is it unique to humans? Behavioral and Brain Sciences, 30(3), 299-351. \url{https://doi.org/10.1017/S0140525X07001975}
\bibitem{ref7} Schacter, D. L., Benoit, R. G., \& Szpunar, K. K. (2017). Episodic future thinking: Mechanisms and functions. Current Opinion in Behavioral Sciences, 17, 41-50. \url{https://doi.org/10.1016/j.cobeha.2017.06.002}
\bibitem{ref8} Schacter, D. L., \& Addis, D. R. (2007). The cognitive neuroscience of constructive memory: Remembering the past and imagining the future. Philosophical Transactions of the Royal Society B, 362(1481), 773-786. \url{https://doi.org/10.1098/rstb.2007.2087}
\bibitem{ref9} Szpunar, K. K., Spreng, R. N., \& Schacter, D. L. (2014). A taxonomy of prospection: Introducing an organizational framework for future-oriented cognition. PNAS, 111(52), 18414-18421. \url{https://doi.org/10.1073/pnas.1417144111}
\bibitem{ref10} Karger, E., Bastani, H., Yueh-Han, C., et al. (2024). ForecastBench: A Dynamic Benchmark of AI Forecasting Capabilities. arXiv:2409.19839. \url{https://arxiv.org/abs/2409.19839}
\bibitem{ref11} Fatemi, B., Kazemi, M., Tsitsulin, A., et al. (2024). Test of Time: A Benchmark for Evaluating LLMs on Temporal Reasoning. arXiv:2406.09170. \url{https://arxiv.org/abs/2406.09170}
\bibitem{ref12} Wu, D., Wang, H., Yu, W., et al. (2025). LongMemEval: Benchmarking Chat Assistants on Long-Term Interactive Memory. ICLR 2025. arXiv:2410.10813. \url{https://arxiv.org/abs/2410.10813}
\bibitem{ref13} Bean, A. M., et al. (2025). Measuring What Matters: Construct Validity in Large Language Model Benchmarks. NeurIPS 2025. arXiv:2511.04703. \url{https://arxiv.org/abs/2511.04703}
\bibitem{ref14} Weidinger, L., Raji, I. D., Wallach, H., et al. (2025). Toward an Evaluation Science for Generative AI Systems. arXiv:2503.05336. \url{https://arxiv.org/abs/2503.05336}
\bibitem{ref15} Wallach, H., Desai, M., Pangakis, N., et al. (2025). Position: Evaluating Generative AI Systems Is a Social Science Measurement Challenge. ICML 2025. arXiv:2502.00561. \url{https://arxiv.org/abs/2502.00561}
\bibitem{ref16} Singh, S., Nan, Y., Wang, A., et al. (2025). The Leaderboard Illusion. arXiv:2504.20879. \url{https://arxiv.org/abs/2504.20879}
\bibitem{ref17} Shridhar, M., Yuan, X., Côté, M.-A., et al. (2021). ALFWorld: Aligning Text and Embodied Environments for Interactive Learning. ICLR 2021. arXiv:2010.03768. \url{https://arxiv.org/abs/2010.03768}
\bibitem{ref18} Côté, M.-A., Kádár, Á., Yuan, X., et al. (2018). TextWorld: A Learning Environment for Text-based Games. arXiv:1806.11532. \url{https://arxiv.org/abs/1806.11532}
\bibitem{ref19} Hausknecht, M., Ammanabrolu, P., Côté, M.-A., \& Yuan, X. (2020). Interactive Fiction Games: A Colossal Adventure (Jericho). AAAI 2020. arXiv:1909.05398. \url{https://arxiv.org/abs/1909.05398}
\bibitem{ref20} Xu, F. F., Song, Y., Li, B., et al. (2024). TheAgentCompany: Benchmarking LLM Agents on Consequential Real World Tasks. arXiv:2412.14161. \url{https://arxiv.org/abs/2412.14161}
\bibitem{ref21} Xie, T., Zhang, D., Chen, J., et al. (2024). OSWorld: Benchmarking Multimodal Agents for Open-Ended Tasks in Real Computer Environments. NeurIPS 2024. arXiv:2404.07972. \url{https://arxiv.org/abs/2404.07972}
\bibitem{ref22} Trivedi, H., Khot, T., Hartmann, M., et al. (2024). AppWorld: A Controllable World of Apps and People for Benchmarking Interactive Coding Agents. ACL 2024. arXiv:2407.18901. \url{https://arxiv.org/abs/2407.18901}
\bibitem{ref23} Yao, S., Shinn, N., Razavi, P., \& Narasimhan, K. (2024). tau-bench: A Benchmark for Tool-Agent-User Interaction in Real-World Domains. arXiv:2406.12045. \url{https://arxiv.org/abs/2406.12045}
\bibitem{ref24} Lu, J., Holleis, T., Zhang, Y., et al. (2024). ToolSandbox: A Stateful, Conversational, Interactive Evaluation Benchmark for LLM Tool Use Capabilities. arXiv:2408.04682. \url{https://arxiv.org/abs/2408.04682}
\bibitem{ref25} Yao, S., Chen, H., Yang, J., \& Narasimhan, K. (2022). WebShop: Towards Scalable Real-World Web Interaction with Grounded Language Agents. NeurIPS 2022. arXiv:2207.01206. \url{https://arxiv.org/abs/2207.01206}
\bibitem{ref26} Deng, X., Gu, Y., Zheng, B., et al. (2023). Mind2Web: Towards a Generalist Agent for the Web. NeurIPS 2023. arXiv:2306.06070. \url{https://arxiv.org/abs/2306.06070}
\bibitem{ref27} Kiela, D., Bartolo, M., Nie, Y., et al. (2021). Dynabench: Rethinking Benchmarking in NLP. NAACL 2021. arXiv:2104.14337. \url{https://arxiv.org/abs/2104.14337}
\bibitem{ref28} Hu, W., et al. (2025). DynaCode: A Dynamic Complexity-Aware Code Benchmark for Evaluating Large Language Models in Code Generation. Findings of ACL 2025. arXiv:2503.10452. \url{https://arxiv.org/abs/2503.10452}
\bibitem{ref29} Zhu, K., Wang, J., Zhao, Q., et al. (2024). Dynamic Evaluation of Large Language Models by Meta Probing Agents. ICML 2024. arXiv:2402.14865. \url{https://arxiv.org/abs/2402.14865}
\bibitem{ref30} Bowman, S. R., \& Dahl, G. E. (2021). What Will it Take to Fix Benchmarking in Natural Language Understanding? NAACL 2021. \url{https://doi.org/10.18653/v1/2021.naacl-main.385}
\bibitem{ref31} Tan, Q., Ng, H. T., \& Bing, L. (2023). TimelineQA: A Benchmark for Question Answering over Timelines. Findings of ACL 2023. \url{https://doi.org/10.18653/v1/2023.findings-acl.6}
\bibitem{ref32} Wei, Y., et al. (2025). TIME: A Multi-level Benchmark for Temporal Reasoning of LLMs in Real-World Scenarios. NeurIPS 2025 Spotlight. arXiv:2505.12891. \url{https://arxiv.org/abs/2505.12891}
\bibitem{ref33} Uddin, M. N., et al. (2025). UnSeenTimeQA: Time-Sensitive Question-Answering Beyond LLMs' Memorization. ACL 2025. arXiv:2407.03525. \url{https://arxiv.org/abs/2407.03525}
\bibitem{ref34} Zou, A., Xiao, T., Jia, R., et al. (2022). Forecasting Future World Events with Neural Networks (AutoCast). NeurIPS 2022. arXiv:2206.15474. \url{https://arxiv.org/abs/2206.15474}
\bibitem{ref35} Halawi, D., Zhang, F., Yueh-Han, C., \& Steinhardt, J. (2024). Approaching Human-Level Forecasting with Language Models. NeurIPS 2024. arXiv:2402.18563. \url{https://arxiv.org/abs/2402.18563}
\bibitem{ref36} Hu, Y., Wang, Y., \& McAuley, J. (2025). Evaluating Memory in LLM Agents via Incremental Multi-Turn Interactions (MemoryAgentBench). arXiv:2507.05257. \url{https://arxiv.org/abs/2507.05257}
\bibitem{ref37} Packer, C., Wooders, S., Lin, K., et al. (2023). MemGPT: Towards LLMs as Operating Systems. arXiv:2310.08560. \url{https://arxiv.org/abs/2310.08560}
\bibitem{ref38} Shinn, N., Cassano, F., Berman, E., et al. (2023). Reflexion: Language Agents with Verbal Reinforcement Learning. NeurIPS 2023. arXiv:2303.11366. \url{https://arxiv.org/abs/2303.11366}
\bibitem{ref39} Park, J. S., O'Brien, J. C., Cai, C. J., et al. (2023). Generative Agents: Interactive Simulacra of Human Behavior. UIST 2023. arXiv:2304.03442. \url{https://arxiv.org/abs/2304.03442}
\bibitem{ref40} Wang, G., Xie, Y., Jiang, Y., et al. (2023). Voyager: An Open-Ended Embodied Agent with Large Language Models. arXiv:2305.16291. \url{https://arxiv.org/abs/2305.16291}
\bibitem{ref41} Ma, C., Zhang, J., Zhu, Z., et al. (2024). AgentBoard: An Analytical Evaluation Board of Multi-Turn LLM Agents. NeurIPS 2024 Oral. arXiv:2401.13178. \url{https://arxiv.org/abs/2401.13178}
\bibitem{ref42} He, P., Dai, Z., He, B., et al. (2025). TRAJECT-Bench: A Trajectory-Aware Benchmark for Evaluating Agentic Tool Use. arXiv:2510.04550. \url{https://arxiv.org/abs/2510.04550}
\bibitem{ref43} Srivastava, A., et al. (2022). Beyond the Imitation Game: Quantifying and Extrapolating the Capabilities of Language Models (BIG-Bench). TMLR. arXiv:2206.04615. \url{https://arxiv.org/abs/2206.04615}
\bibitem{ref44} Gururangan, S., Swayamdipta, S., Levy, O., et al. (2018). Annotation Artifacts in Natural Language Inference Data. NAACL 2018. arXiv:1803.02324. \url{https://arxiv.org/abs/1803.02324}
\bibitem{ref45} Yao, S., Zhao, J., Yu, D., et al. (2023). ReAct: Synergizing Reasoning and Acting in Language Models. ICLR 2023. arXiv:2210.03629. \url{https://arxiv.org/abs/2210.03629}
\bibitem{ref46} Fu, D., et al. (2025). PreAct: Prediction Enhances Agent's Planning Ability. COLING 2025. \url{https://aclanthology.org/2025.coling-main.1/}
\bibitem{ref47} Majumder, B. P., Mishra, B. D., Jansen, P., et al. (2023). CLIN: A Continually Learning Language Agent for Rapid Task Adaptation and Generalization. arXiv:2310.10134. \url{https://arxiv.org/abs/2310.10134}
\bibitem{ref48} Wang, Z. Z., Mao, J., Fried, D., \& Neubig, G. (2025). Agent Workflow Memory. ICML 2025. arXiv:2409.07429. \url{https://arxiv.org/abs/2409.07429}
\bibitem{ref49} Madaan, L., Singh, A. K., Schaeffer, R., et al. (2024). Quantifying Variance in Evaluation Benchmarks. arXiv:2406.10229. \url{https://arxiv.org/abs/2406.10229}
\bibitem{ref50} Blackwell, R. E., Barry, J., \& Cohn, A. G. (2024). Towards Reproducible LLM Evaluation: Quantifying Uncertainty in LLM Benchmark Scores. arXiv:2410.03492. \url{https://arxiv.org/abs/2410.03492}
\bibitem{ref51} Zheng, L., Chiang, W.-L., Sheng, Y., et al. (2023). Judging LLM-as-a-Judge with MT-Bench and Chatbot Arena. NeurIPS 2023. arXiv:2306.05685. \url{https://arxiv.org/abs/2306.05685}
\bibitem{ref52} Lee, Y., et al. (2025). CheckEval: A Reliable LLM-as-a-Judge Framework for Evaluating Text Generation Using Checklists. EMNLP 2025. \url{https://doi.org/10.18653/v1/2025.emnlp-main.796}
\bibitem{ref53} Mialon, G., Fourrier, C., Swift, C., et al. (2023). GAIA: A Benchmark for General AI Assistants. arXiv:2311.12983. \url{https://arxiv.org/abs/2311.12983}
\bibitem{ref54} Rein, D., et al. (2025). HCAST: Human-Calibrated Autonomy Software Tasks. arXiv:2503.17354. \url{https://arxiv.org/abs/2503.17354}
\bibitem{ref55} Wijk, H., Lin, T., Becker, J., et al. (2024). RE-Bench: Evaluating Frontier AI R\&D Capabilities of Language Model Agents against Human Experts. arXiv:2411.15114. \url{https://arxiv.org/abs/2411.15114}
\bibitem{ref56} Wester, J., et al. (2022). A Framework for Rigorous Evaluation of Human Performance in Human and Machine Learning Comparison Studies. Scientific Reports, 12, 5400. \url{https://doi.org/10.1038/s41598-022-08831-8}
\bibitem{ref57} Lee, M., Srivastava, M., Hardy, A., et al. (2023). Evaluating Human-Language Model Interaction (HALIE). TMLR. arXiv:2212.09746. \url{https://arxiv.org/abs/2212.09746}
\bibitem{ref58} Lal, Y. K., Chambers, N., Mooney, R., \& Balasubramanian, N. (2021). TellMeWhy: A Dataset for Answering Why-Questions in Narratives. Findings of ACL-IJCNLP 2021. \url{https://doi.org/10.18653/v1/2021.findings-acl.53}
\bibitem{ref59} National Transportation Safety Board. (n.d.). CAROL Query: Accident Data. \url{https://data.ntsb.gov/carol-main-public/}
\bibitem{ref60} Xu, Y., Wang, D., Yu, M., et al. (2022). Fantastic Questions and Where to Find Them: FairytaleQA --- An Authentic Dataset for Narrative Comprehension. ACL 2022. \url{https://doi.org/10.18653/v1/2022.acl-long.34}
\end{thebibliography}
